\documentclass[UTF8,a4paper,11pt]{ctexart}
\usepackage{amsmath,amssymb,bm,geometry,fancyhdr,booktabs,array}
\geometry{left=22mm,right=22mm,top=22mm,bottom=22mm}
\pagestyle{fancy}\fancyhf{}
\lhead{过程控制工程·第三周作业标准解答}
\rhead{PID 整定：ZN 与 Lambda}
\cfoot{\thepage}
\setlength{\parindent}{2em}
\setlength{\parskip}{0.35em}
\begin{document}
\begin{center}
{\LARGE\bfseries 过程控制工程·第三周作业标准解答}\\[4pt]
{\large 管式换热器开环阶跃试验与 PID 参数整定}
\end{center}

\section*{1. 已知数据与控制对象}
蒸汽压力由 $18\,\mathrm{psig}$ 阶跃到 $20\,\mathrm{psig}$，故
\[
\Delta P_s=2\,\mathrm{psi}.
\]
温度变送器输出由约 $12.0\,\mathrm{mA}$ 上升到约 $17.0\,\mathrm{mA}$，故
\[
\Delta T_{2m}=5.0\,\mathrm{mA}.
\]
因此换热器与测量环节从蒸汽压力到温度测量信号的稳态增益为
\[
K_p=\frac{\Delta T_{2m}}{\Delta P_s}=\frac{5}{2}=2.5\;\mathrm{mA/psi}.
\]
给定
\[
K_{IP}=0.75\;\mathrm{psi/mA},\qquad K_v=0.9\;\mathrm{psi/psi},
\]
则从控制器输出到温度测量信号的总稳态增益为
\[
\boxed{K=K_{IP}K_vK_p=0.75\times0.9\times2.5=1.6875}.
\]

\section*{2. 由阶跃响应辨识一阶惯性纯滞后模型}
采用课程小测中使用的 $28.3\%$--$63.2\%$ 两点法。初值与终值为
\[
y_0=12.0,\qquad y_\infty\approx17.0,\qquad \Delta y=5.0.
\]
因此
\[
y_{28.3\%}=12+0.283\times5=13.415\,\mathrm{mA},
\]
\[
y_{63.2\%}=12+0.632\times5=15.160\,\mathrm{mA}.
\]
由表中线性插值得
\[
t_{28.3\%}\approx3+\frac{13.415-13.1}{14.0-13.1}=3.35\,\mathrm{min},
\]
\[
t_{63.2\%}\approx5+\frac{15.160-14.8}{15.4-14.8}=5.60\,\mathrm{min}.
\]
按课程两点公式
\[
T=1.5\left(t_{63.2\%}-t_{28.3\%}\right)
=1.5(5.60-3.35)=3.375\,\mathrm{min},
\]
\[
\tau=t_{63.2\%}-T=5.60-3.375=2.225\,\mathrm{min}.
\]
故
\[
\boxed{G(s)=\frac{1.6875e^{-2.225s}}{3.375s+1}}.
\]

\section*{3. Ziegler--Nichols 法}
按课程小测使用的 ZN 形式
\[
K_c=\frac{1}{K}\frac{T}{\tau},\qquad
T_i=2\tau,\qquad
T_d=0.5\tau.
\]
代入得
\[
K_c=0.899,\qquad
T_i=4.450\,\mathrm{min},\qquad
T_d=1.113\,\mathrm{min}.
\]
因此
\[
\boxed{K_c^{ZN}=0.899,\qquad T_i^{ZN}=4.45\,\mathrm{min},\qquad T_d^{ZN}=1.11\,\mathrm{min}}.
\]
若采用理想 PID
\[
G_c(s)=K_c\left(1+\frac{1}{T_i s}+T_d s\right),
\]
则
\[
\boxed{G_c^{ZN}(s)=0.899\left(1+\frac{1}{4.45s}+1.113s\right)}.
\]

\section*{4. Lambda 法}
按课程小测给出的规则取
\[
\lambda=0.2T=0.675\,\mathrm{min}.
\]
整定公式为
\[
K_c=\frac{1}{K}\frac{T}{\tau+\lambda},\qquad
T_i=T,\qquad
T_d=\frac{\tau}{2}.
\]
故
\[
K_c=0.690,\qquad
T_i=3.375\,\mathrm{min},\qquad
T_d=1.113\,\mathrm{min}.
\]
因此
\[
\boxed{K_c^{\lambda}=0.690,\qquad T_i^{\lambda}=3.375\,\mathrm{min},\qquad T_d^{\lambda}=1.113\,\mathrm{min}}.
\]
相应控制器为
\[
\boxed{G_c^{\lambda}(s)=0.690\left(1+\frac{1}{3.375s}+1.113s\right)}.
\]

\section*{5. 汇总}
\begin{center}
\begin{tabular}{lccc}
\toprule
方法 & $K_c$ & $T_i$/min & $T_d$/min\\
\midrule
ZN & 0.899 & 4.45 & 1.11\\
Lambda ($\lambda=0.2T$) & 0.690 & 3.375 & 1.11\\
\bottomrule
\end{tabular}
\end{center}

\paragraph{说明}
若课堂对 Lambda 参数 $\lambda$ 指定其他取值，保留已辨识的 $K,T,\tau$，重新代入 Lambda 公式即可。
\end{document}
